In Section~\ref{sec:design}, we discussed how to create designs in Orbiter. 
In this section, we wish to compute properties of designs related to the automorphism group.
The first step is to 
export the incidence matrix of the design to file. 
After that, we compute the canonical form of the design, 
which allows us to determine many properties.
The following example computes the properties of $\PG(2,3)$:

\sourcecode{design_PG_2_3_canonical}






\bigskip


The command 

\sourcecode{wreath_product_designs_n4_k2_c}

computes the automorphism group of the design on 8 points created in Section~\ref{sec:design}.
The group is $\Sym(4) \wr \Sym(2)$.
The command 

\sourcecode{wreath_product_designs_n8_k6_c}

computes the automorphism group of the design on 16 points created in Section~\ref{sec:design}.
The group is $\Sym(8) \wr \Sym(2)$.


\bigskip



In Section~\ref{sec:large:sets}, the large sets of $\AG(2,3)$ were constructed. 
However, the isomorph classsification of the created objects is still outstanding. 
We will do this final step now.
The isomorphism classification is performed using the Nauty interface. 
A list of combinatorial objects is created, 
and the  \verb'-canonical_form' command is applied as activity.
This will produce a list of pairwise non-isomorphic designs.
The size of this list is the number of isomorphism types of large sets of $\AG(2,3).$

\sourcecode{LS_AG_2_3_solutions_classify}

It turns out that there are exactly two isomorphism types, 
with automorphism groups of order 54 and 42, respectively.

\bigskip

